% GENERATED FILE. Edit canonical lessons, then run npm run generate:books.
% canonical-source-hash: fnv1a64:5cc0b729b5b2ccd4
% canonical-lessons: HI-C11-kaalaa-safed, HI-C11-laal-niila

\chapter{Four Core Colours}
\label{ch:hi-core-colours}

This chapter is generated from the canonical micro-lessons used by Language Ladder.
Each section stays independently resumable and preserves the authored prerequisite order.

\section[\hi{काला} \hi{सफ़ेद}]{\hi{काला}, \hi{सफ़ेद} (kālā, safed) — black and white}

\label{lesson:HI-C11-kaalaa-safed}



\begin{quote}
\textbf{Warm-up.} {[}PAUSE 2s{]} Hindi's word for black hides something bigger inside it: the very same word Sanskrit uses for \textbf{time itself} — and for death.
\end{quote}

\subsection*{You'll want to know — \hi{काला} — black, time, and death, tangled together}
\textbf{\hi{काला}} (\emph{kālā}) = "\textbf{black}," from Sanskrit \textbf{\hi{काल}} (\emph{kāla}), a word that covers both "\textbf{black/dark}" \textbf{and} "\textbf{time}." Sanskrit's own tradition ties the two together poetically — \textbf{time devours everything}, and that devouring is imagined as \emph{dark} — and the same word names the fierce goddess \textbf{Kālī} ("the black/dark one") and stands as an epithet of \textbf{Yama}, the god of death. Be honest, though: whether "black" and "time" are truly \textbf{one and the same root}, or two separate ancient words that Sanskrit poetry and theology later \textbf{fused} into a single idea, is a real, still-debated question among linguists — the cultural connection is certain; the precise etymology underneath it is not.

\subsection*{You'll want to know — \hi{सफ़ेद} — a Persian loan}
\textbf{\hi{सफ़ेद}} (\emph{safed}) = "\textbf{white}" — borrowed from \textbf{Persian} (\ar{سفید} \emph{safēd}), the same kind of Persian-into-Hindi borrowing you've already met in \textbf{\hi{माफ़}} (\emph{māf}, "forgiven," from the "sorry" lesson). Notice the \textbf{\hi{नुक़्ता}} (\emph{nuqtā}, the dot under \hi{फ}): it turns \emph{pha} into \textbf{fa}, the tell-tale mark of a sound that arrived from Persian or Arabic.

\subsection*{Guided Practice}
{[}PAUSE 1s{]}

\begin{itemize}
\raggedright
  \item {[}YOU SAY: "kāla" — time/black/death, then "kālā" — black{]}
  \item {[}YOU SAY: "safed" — white, a Persian loan{]}
  \item {[}YOU SAY: the dot — \hi{नुक़्ता} under \hi{फ} turns pha into fa{]}
\end{itemize}

\begin{tcolorbox}[breakable,colback=teal!4,colframe=teal!35!black,title={Wrap-up recall}]
{[}PAUSE 3s{]} What does \textbf{\hi{काला}}'s root \textbf{\hi{काल}} (\emph{kāla}) also mean, besides "black"? (\textbf{"Time"} — and by extension, \textbf{death}; the same word names the goddess \textbf{Kālī} and Yama's epithet \textbf{Kāla}.) Is it settled that "black" and "time" are the same root? (\textbf{No} — the cultural link is real and ancient, but linguists still debate whether it's one true root or two words later fused.) Where does \textbf{\hi{सफ़ेद}} come from? (\textbf{Persian} \emph{safēd} — the same kind of loan as \emph{māf}, "sorry.")
\end{tcolorbox}
\section[\hi{लाल} \hi{नीला}]{\hi{लाल}, \hi{नीला} (lāl, nīlā) — a gemstone's name, and India's dye}

\label{lesson:HI-C11-laal-niila}



\begin{quote}
\textbf{Warm-up.} {[}PAUSE 2s{]} One more color pair, and one more Persian loan — plus a Sanskrit word connected to why English itself says "\textbf{indigo}."
\end{quote}

\subsection*{You'll want to know — \hi{लाल} — named for a gem, not a color}
\textbf{\hi{लाल}} (\emph{lāl}) = "\textbf{red}" — borrowed from Persian \textbf{\ar{لعل}} (properly \emph{laʿl}, with a middle \emph{ʿayn} the Hindi/Persian pronunciation \emph{lāl} smooths over), which originally named a specific \textbf{deep-red gemstone} (a ruby or spinel), and only later widened to mean the color itself. Hindi's everyday word for "red" started life as a jeweler's word.

\subsection*{You'll want to know — \hi{नीला} — Sanskrit's dark blue, and the world's word for indigo}
\textbf{\hi{नीला}} (\emph{nīlā}) = "\textbf{blue}," from Sanskrit \textbf{\hi{नील}} (\emph{nīla}), "dark blue" — also the name of the \textbf{indigo plant} and its dye. Here's the honest, careful version of a famous connection: English \textbf{"indigo"} does \textbf{not} come from \emph{nīla} by sound change. It comes from \textbf{Greek} \emph{indikón}, simply meaning "\textbf{the Indian {[}thing{]}}" — because Greek and Roman traders imported the blue dye \textbf{from India}. So \emph{nīlā} and "indigo" aren't the same word evolving — they're two \textbf{separate} words for the same dye, one Indian (\emph{nīla}, describing the color) and one foreign (\emph{indikón}, describing the dye's \emph{origin}). The connection is real, but it's a connection of \textbf{history and trade}, not of sound.

\subsection*{Guided Practice}
{[}PAUSE 1s{]}

\begin{itemize}
\raggedright
  \item {[}YOU SAY: "lāl" — red, originally a gemstone's name{]}
  \item {[}YOU SAY: "nīlā" — blue, from Sanskrit nīla{]}
  \item {[}YOU SAY: the honest distinction — nīla and "indigo" are two different words linked by trade, not by sound{]}
\end{itemize}

\begin{tcolorbox}[breakable,colback=teal!4,colframe=teal!35!black,title={Wrap-up recall}]
{[}PAUSE 3s{]} What did \textbf{\hi{लाल}} originally name, before it meant "red"? (\textbf{A deep red gemstone}, borrowed from Persian.) Does English "\textbf{indigo}" come directly from Sanskrit \textbf{\hi{नील}} \emph{nīla}? (\textbf{No} — indigo is from Greek \emph{indikón}, "the Indian {[}thing{]}," describing where the dye came from; \emph{nīla} and "indigo" are two separate words connected by history, not by sound change.)
\end{tcolorbox}
